\documentclass[12pt,a4paper]{article}

% 宏包设置
\usepackage{ctex}
\usepackage{amsmath, amssymb}
\usepackage{geometry}
\geometry{left=2cm, right=2cm, top=2.5cm, bottom=2.5cm}
\usepackage{enumitem}
\usepackage{fancyhdr}

% 页眉页脚设置


\title{\textbf{微积分训练}}
\author{}
\date{}

\begin{document}

\maketitle
\thispagestyle{fancy}

\section*{微积分 I}
\textbf{说明：} 请选择出每道题的正确答案。\\
\textbf{分值：} 每题 1 分。

\begin{enumerate}[label=\arabic*), itemsep=15pt]

    % Q1: 绝对值极限
    \item 计算极限：$\lim_{x\rightarrow 4^-} \frac{16-x^2}{|x-4|} =$
    \begin{enumerate}[label=\alph*), itemsep=0pt]
        \item 0 \quad b) 8 \quad c) -8 \quad d) 16 \quad e) $\infty$
    \end{enumerate}

    % Q2: 分段函数极限
    \item 若分段函数 $f(x) = \begin{cases} 3x^2 - a, & \text{如果 } x \ge 2 \\ 2x + 5, & \text{如果 } x < 2 \end{cases}$ 在 $x \to 2$ 时的极限 $\lim_{x\rightarrow 2} f(x)$ 存在，则常数 $a$ 的值为：
    \begin{enumerate}[label=\alph*), itemsep=0pt]
        \item 3 \quad b) 9 \quad c) 12 \quad d) -3 \quad e) 不存在
    \end{enumerate}

    % Q3: 导数的极限定义
    \item 已知函数 $f(x) = \frac{2}{x} + x^2$，则极限 $\lim_{h\rightarrow 0} \frac{f(3+h)-f(3)}{h} =$
    \begin{enumerate}[label=\alph*), itemsep=0pt]
        \item $\frac{52}{9}$ \quad b) $\frac{56}{9}$ \quad c) $6$ \quad d) $\frac{16}{3}$ \quad e) $\frac{50}{9}$
    \end{enumerate}

    % Q4: 导数与二阶导数的切线斜率比值
    \item 设 $f(x) = x^3 - 2x^2$。函数 $f''(x)$ 图像在 $x=2$ 处的切线斜率，与函数 $f(x)$ 图像在 $x=2$ 处的切线斜率，这两者的比值是：
    \begin{enumerate}[label=\alph*), itemsep=0pt]
        \item $\frac{3}{2}$ \quad b) $\frac{4}{3}$ \quad c) $\frac{2}{3}$ \quad d) $\frac{3}{4}$ \quad e) $1$
    \end{enumerate}

    % Q5: 函数凹凸性与增减性描述
    \item 对于形如 $f(x) = (cx-d)^4$ 的函数，其中 $c > d > 0$。当 $x=0$ 时，下列哪项描述必定正确？\\
    I. 曲线正在以递增的速率增加。\\
    II. 曲线正在以递增的速率减少。\\
    III. 曲线正在以递减的速率减少。
    \begin{enumerate}[label=\alph*), itemsep=0pt]
        \item 仅 I \quad b) 仅 II \quad c) 仅 III \quad d) I 和 II \quad e) II 和 III
    \end{enumerate}

    % Q6: 线性近似
    \item 已知 $f(4) = 3$ 且 $f'(4) = -\frac{1}{2}$，则 $f(3.96)$ 的近似值为：
    \begin{enumerate}[label=\alph*), itemsep=0pt]
        \item $2.98$ \quad b) $3.02$ \quad c) $3.04$ \quad d) $2.96$ \quad e) $3.08$
    \end{enumerate}

    % Q7: 切线问题
    \item 直线 $y = 6x - 5$ 是曲线 $y = 2x^2 - 2x + 3$ 在哪一点的切线？
    \begin{enumerate}[label=\alph*), itemsep=0pt]
        \item $(1, 1)$ \quad b) $(2, 7)$ \quad c) $(3, 13)$ \quad d) $(0, -5)$ \quad e) $(2, 5)$
    \end{enumerate}

    % Q8: 均值定理
    \item 针对函数 $f(x) = \sqrt{x+1}$ 在闭区间 $[0, 3]$ 上，求满足中值定理（Mean Value Theorem）的常数 $c$ 的值：
    \begin{enumerate}[label=\alph*), itemsep=0pt]
        \item $\frac{3}{4}$ \quad b) $\frac{5}{4}$ \quad c) $\frac{4}{3}$ \quad d) $1$ \quad e) $\frac{5}{2}$
    \end{enumerate}

    % Q9: 极限计算进阶
    \item 计算极限：$\lim_{x \rightarrow \infty} \frac{\sqrt{9x^2 + 2}}{4x - 3} =$
    \begin{enumerate}[label=\alph*), itemsep=0pt]
        \item $\frac{9}{4}$ \quad b) $\frac{3}{4}$ \quad c) $-\frac{3}{4}$ \quad d) $0$ \quad e) $\infty$
    \end{enumerate}

    % Q10: 导数的变体极限
    \item 计算极限：$\lim_{h \rightarrow 0} \frac{\cos(\frac{\pi}{3} + h) - \frac{1}{2}}{h} =$
    \begin{enumerate}[label=\alph*), itemsep=0pt]
        \item $\frac{\sqrt{3}}{2}$ \quad b) $-\frac{\sqrt{3}}{2}$ \quad c) $\frac{1}{2}$ \quad d) $-\frac{1}{2}$ \quad e) $0$
    \end{enumerate}

    % Q11: 隐函数求导基础
    \item 已知隐函数方程 $x^2 y + y^2 x = 6$，求在点 $(1, 2)$ 处的导数 $\frac{dy}{dx}$ 的值：
    \begin{enumerate}[label=\alph*), itemsep=0pt]
        \item $-\frac{4}{5}$ \quad b) $-1$ \quad c) $-\frac{8}{5}$ \quad d) $-\frac{5}{4}$ \quad e) $1$
    \end{enumerate}

    % Q12: 极值判断
    \item 函数 $f(x) = x^4 - 4x^3$ 有几个相对极值点（极大值或极小值）？
    \begin{enumerate}[label=\alph*), itemsep=0pt]
        \item 0 \quad b) 1 \quad c) 2 \quad d) 3 \quad e) 4
    \end{enumerate}

    % Q13: 中值定理代数应用
    \item 已知 $f(x) = x^3 + x - 1$，根据中值定理，在区间 $[0, 2]$ 内必然存在一点 $c$，使得 $f'(c) =$
    \begin{enumerate}[label=\alph*), itemsep=0pt]
        \item 1 \quad b) 3 \quad c) 5 \quad d) 7 \quad e) 9
    \end{enumerate}

    % Q14: 高阶导数
    \item 如果 $f(x) = x \ln x$，求二阶导数在 $x=e$ 处的值 $f''(e) =$
    \begin{enumerate}[label=\alph*), itemsep=0pt]
        \item $1$ \quad b) $\frac{1}{e}$ \quad c) $e$ \quad d) $0$ \quad e) $\frac{2}{e}$
    \end{enumerate}

    % Q15: 切线截距
    \item 曲线 $y = \frac{1}{x^2}$ 在 $x = 1$ 处的切线，其在 $y$ 轴上的截距为：
    \begin{enumerate}[label=\alph*), itemsep=0pt]
        \item 1 \quad b) 2 \quad c) 3 \quad d) -1 \quad e) -2
    \end{enumerate}

\end{enumerate}

\newpage
\section*{微积分 II}
\textbf{说明：} 请选择出每道题的正确答案。\\
\textbf{分值：} 每题 1 分。

\begin{enumerate}[label=\arabic*), itemsep=15pt, start=16]

    % Q16: 对称差商极限
    \item 极限 $\lim_{h\rightarrow 0} \frac{f(3+2h)-f(3-2h)}{h} =$
    \begin{enumerate}[label=\alph*), itemsep=0pt]
        \item $f'(3)$ \quad b) $2f'(3)$ \quad c) $4f'(3)$ \quad d) $\frac{1}{2}f'(3)$ \quad e) 0
    \end{enumerate}

    % Q17: 水平渐近线
    \item 函数 $f(x) = \frac{(2x^2 - 1)(3x^3 + 2)}{(x^2 + 5)(x^3 - 4x)}$ 的水平渐近线方程是：
    \begin{enumerate}[label=\alph*), itemsep=0pt]
        \item $y = 6$ \quad b) $y = 5$ \quad c) $y = 2$ \quad d) $y = \frac{2}{5}$ \quad e) $y = 0$
    \end{enumerate}

    % Q18: 运动学 (加速度)
    \item 一个质点沿 x 轴运动，其在时间 $t$ 的位置函数由 $x = \frac{2t-1}{t+2}$ 给出。该质点在 $t = -1$ 时的加速度是：
    \begin{enumerate}[label=\alph*), itemsep=0pt]
        \item $5$ \quad b) $10$ \quad c) $-10$ \quad d) $-5$ \quad e) $0$
    \end{enumerate}

    % Q19: 复合求导比率
    \item 已知 $y = 3x^4 - 2$，求导数比率 $\frac{dy}{d(x^2+1)} =$
    \begin{enumerate}[label=\alph*), itemsep=0pt]
        \item $6x^2$ \quad b) $12x^2$ \quad c) $6x$ \quad d) $3x^2$ \quad e) $12x$
    \end{enumerate}

    % Q20: 隐函数二阶导数
    \item 对于隐函数 $x^2 - y^2 = 9$，其二阶导数 $\frac{d^2y}{dx^2}$（用含 $y$ 的表达式表示）为：
    \begin{enumerate}[label=\alph*), itemsep=0pt]
        \item $\frac{x}{y^2}$ \quad b) $-\frac{9}{y^3}$ \quad c) $\frac{9}{y^3}$ \quad d) $-\frac{x}{y^3}$ \quad e) $\frac{y^2-x^2}{y^3}$
    \end{enumerate}

    % Q21: 相关变化率 (动态表达式)
    \item 已知 $z = \frac{x^2 - y}{x}$。如果 $x$ 正以 $1 \text{ 单位/秒}$ 的速率增加，且 $y$ 正以 $4 \text{ 单位/秒}$ 的速率增加。当 $x = 2$ 且 $y = 4$ 时，$z$ 增加的速率是多少？
    \begin{enumerate}[label=\alph*), itemsep=0pt]
        \item $-1$ \quad b) $0$ \quad c) $1$ \quad d) $2$ \quad e) $\frac{1}{2}$
    \end{enumerate}

    % Q22: 最优化 (几何)
    \item 将一个矩形内接于半径为 4 的半圆中。这个矩形可能的最大面积是多少？
    \begin{enumerate}[label=\alph*), itemsep=0pt]
        \item 8 \quad b) $8\sqrt{2}$ \quad c) 16 \quad d) $16\sqrt{2}$ \quad e) 32
    \end{enumerate}

    % Q23: 运动学 (重力与抛体)
    \item 一枚模型火箭从地面发射，其净加速度为 $a = 10 - 2t$。在第 5 秒结束时，火箭发动机关闭，此后火箭仅受重力加速度影响（$a = -32 \text{ 英尺/秒}^2$）。假设初始速度 $v(0)=0$，求在 $t=6$ 秒时火箭的速度：
    \begin{enumerate}[label=\alph*), itemsep=0pt]
        \item $-7$ \quad b) $25$ \quad c) $-32$ \quad d) $-12$ \quad e) $0$
    \end{enumerate}

    % Q24: 积分平均值 (预热)
    \item 若函数 $y = 3x^2$ 在区间 $[1, a]$ （其中 $a > 1$）上的平均值为 $13$，则常数 $a$ 的值为：
    \begin{enumerate}[label=\alph*), itemsep=0pt]
        \item 2 \quad b) 3 \quad c) 4 \quad d) 5 \quad e) 6
    \end{enumerate}

    % Q25: 反常/高阶导数概念
    \item 设 $f^{(n)}(x)$ 表示函数 $f(x) = \sin(2x)$ 的第 $n$ 阶导数，求第 10 阶导数 $f^{(10)}(x) =$
    \begin{enumerate}[label=\alph*), itemsep=0pt]
        \item $2^{10} \sin(2x)$ \quad b) $-2^{10} \sin(2x)$ \quad c) $2^{10} \cos(2x)$ \quad d) $-2^{10} \cos(2x)$ \quad e) $2 \sin(2x)$
    \end{enumerate}

    % Q26: 变上限积分导数
    \item 求导：$\frac{d}{dx} \int_{x}^{3x} \sqrt{t^2 + 1} \, dt =$
    \begin{enumerate}[label=\alph*), itemsep=0pt]
        \item $\sqrt{9x^2+1} - \sqrt{x^2+1}$ 
        \item $3\sqrt{9x^2+1} - \sqrt{x^2+1}$ 
        \item $3\sqrt{x^2+1}$
        \item $\sqrt{8x^2}$
        \item $3x\sqrt{9x^2+1}$
    \end{enumerate}

    % Q27: 洛必达法则进阶
    \item 计算极限：$\lim_{x\rightarrow 0} \frac{e^x - x - 1}{x^2} =$
    \begin{enumerate}[label=\alph*), itemsep=0pt]
        \item 0 \quad b) 1 \quad c) $\frac{1}{2}$ \quad d) $\infty$ \quad e) 不存在
    \end{enumerate}

    % Q28: 极坐标面积思想 / 积分
    \item 由曲线 $y = x^2 - 4$ 与 x 轴所围成的封闭区域面积是：
    \begin{enumerate}[label=\alph*), itemsep=0pt]
        \item $\frac{32}{3}$ \quad b) $-\frac{32}{3}$ \quad c) $\frac{16}{3}$ \quad d) $\frac{64}{3}$ \quad e) 0
    \end{enumerate}

    % Q29: 微分方程初步
    \item 大气压强 $P$ 随海拔高度 $h$ 的变化率与当前压强 $P$ 成正比。下列哪一个微分方程描述了这种关系？
    \begin{enumerate}[label=\alph*), itemsep=0pt]
        \item $\frac{dP}{dh} = k$ \quad b) $\frac{dP}{dh} = kh$ \quad c) $\frac{dP}{dh} = kP$ \quad d) $\frac{dh}{dP} = kP$ \quad e) $\frac{dP}{dh} = \frac{k}{P}$
    \end{enumerate}

    % Q30: 旋转体体积基础
    \item 由曲线 $y = \sqrt{x}$，直线 $x=4$ 以及 x 轴围成的区域，绕 x 轴旋转一周所生成的旋转体体积是：
    \begin{enumerate}[label=\alph*), itemsep=0pt]
        \item $8\pi$ \quad b) $16\pi$ \quad c) $\frac{16\pi}{3}$ \quad d) $4\pi$ \quad e) $2\pi$
    \end{enumerate}

\end{enumerate}



\end{document}